% SOURCE: Chinese PDF viewer pp. 84--86 (until heading of Section 6).
\section{Orbital Angular Momentum}
\subsection*{1) Eigenvalues and eigenstates}
The classical orbital angular momentum and its quantum operator are
\[\vect L=\vect r\times\vect p\quad\longrightarrow\quad
\op{\vect L}=\op{\vect r}\times\op{\vect p},\qquad
\op L_i=\epsilon_{ijk}\op x_j\op p_k\quad (i,j,k=1,2,3).\]
From the fundamental commutation relations
\[[\op x_i,\op x_j]=0,\qquad[\op p_i,\op p_j]=0,\qquad
[\op x_i,\op p_j]=\ii\hbar\delta_{ij},\]
one can show that
\[[\op L_i,\op L_j]=\ii\hbar\epsilon_{ijk}\op L_k.\]
Thus the orbital angular-momentum operators satisfy the general definition of angular momentum.

Because this is orbital angular momentum, we calculate eigenvalues and eigenstates in the coordinate representation:
\[\op{\vect r}=\vect r,\qquad\op{\vect p}=-\ii\hbar\vect\nabla,\]
\[\op{\vect L}=-\ii\hbar\vect r\times\vect\nabla,\qquad
\op L_i=-\ii\hbar\epsilon_{ijk}x_j\frac\partial{\partial x_k}.\]
It is convenient to consider rotations in spherical coordinates, $(x,y,z)\to(r,\theta,\varphi)$. The operators are
\[\begin{aligned}
\op L_x&=\ii\hbar\left(\sin\varphi\frac\partial{\partial\theta}
+\cot\theta\cos\varphi\frac\partial{\partial\varphi}\right),\\
\op L_y&=-\ii\hbar\left(\cos\varphi\frac\partial{\partial\theta}
-\cot\theta\sin\varphi\frac\partial{\partial\varphi}\right),\\
\op L_z&=-\ii\hbar\frac\partial{\partial\varphi},
\end{aligned}\]
\[\op{\vect L}^{\,2}=-\hbar^2\left[
\frac1{\sin\theta}\frac\partial{\partial\theta}
\left(\sin\theta\frac\partial{\partial\theta}\right)
+\frac1{\sin^2\theta}\frac{\partial^2}{\partial\varphi^2}\right].\]
This shows that $\op{\vect L}$ depends only on the angles $\theta$ and $\varphi$.

First consider the eigenvalue equation of $\op L_z$. As $\op L_z$ does not depend on $\theta$,
\[\op L_z\Phi(\varphi)=L_z\Phi(\varphi),\qquad
-\ii\hbar\frac{\dd\Phi}{\dd\varphi}=L_z\Phi(\varphi).\]
The solution is
\[\Phi(\varphi)=A\exp(\ii L_z\varphi/\hbar).\]
The single-valuedness condition
\[\Phi(\varphi)=\Phi(\varphi+2\pi)\]
implies $L_z=m\hbar$, where $m=0,\pm1,\ldots,\pm\infty$, and $\Phi(\varphi)=Ae^{\ii m\varphi}$. Normalize:
\[\int_0^{2\pi}|\Phi(\varphi)|^2\dd\varphi=1
\quad\Longrightarrow\quad A=\frac1{\sqrt{2\pi}}.\]
Thus
\[\boxed{\Phi(\varphi)=\frac{e^{\ii m\varphi}}{\sqrt{2\pi}}.}\]

Next consider the eigenvalue equation of $\op{\vect L}^{\,2}$:
\[\op{\vect L}^{\,2}Y(\theta,\varphi)=L^2Y(\theta,\varphi).\]
This is
\[\left[
\frac1{\sin\theta}\frac\partial{\partial\theta}
\left(\sin\theta\frac\partial{\partial\theta}\right)
+\frac1{\sin^2\theta}\frac{\partial^2}{\partial\varphi^2}
+\frac{L^2}{\hbar^2}\right]Y_{lm}(\theta,\varphi)=0.\]
Using separation of variables and the single-valuedness and finiteness conditions, one finds
\[L^2=l(l+1)\hbar^2,\qquad l=0,1,2,\ldots.\]
The eigenstates are spherical harmonics
\[Y_{lm}(\theta,\varphi)=N_{lm}P_l^m(\cos\theta)e^{\ii m\varphi},\]
where $P_l^m(x)$ is the associated Legendre polynomial. The magnetic quantum number is
\[m=0,\pm1,\pm2,\ldots,\pm l,\]
with $2l+1$ values. The eigenvalue $L^2$ depends only on $l$ whereas the eigenstates also depend on $m$, so the degeneracy is $g=2l+1$.

The normalization condition is
\[\int_0^{2\pi}\!\dd\varphi\int_0^\pi\!\sin\theta\dd\theta\,
|Y_{lm}(\theta,\varphi)|^2=1,\]
which gives the normalization factor
\[N_{lm}=\sqrt{\frac{(l-m)!(2l+1)}{(l+m)!4\pi}}.\]
Because the $\varphi$-dependent factor in $Y_{lm}$ is an eigenstate of $\op L_z$, the spherical harmonics are common eigenstates of $\op{\vect L}^{\,2}$ and $\op L_z$:
\[\boxed{\begin{aligned}
\op{\vect L}^{\,2}Y_{lm}(\theta,\varphi)&=l(l+1)\hbar^2Y_{lm}(\theta,\varphi),\\
\op L_zY_{lm}(\theta,\varphi)&=m\hbar Y_{lm}(\theta,\varphi),
\end{aligned}}\]
where $l=0,1,2,\ldots$ and $m=0,\pm1,\ldots,\pm l$. The component of angular momentum is constrained by its total magnitude, so $m$ cannot extend to $\pm\infty$ for fixed $l$.

This shows that orbital angular momentum realizes the general angular-momentum algebra in coordinate space. It is only one form of angular momentum: $l$ is restricted to zero and positive integers and cannot be a positive half-integer.

\subsection*{2) Spatial rotation operator}
The orbital-angular-momentum eigenstate is the projection of $\ket{l,m}$ onto a radial position basis state $\ket{\vect e_r(\theta,\varphi)}$:
\[Y_{lm}(\theta,\varphi)=\braket{\vect e_r(\theta,\varphi)}{l,m}.\]
The unit vector $\vect e_r(\theta,\varphi)$ can be obtained from $\vect e_z(0,\varphi)$ by two rotations, first by $\theta$ about the $y$ axis and then by $\varphi$ about the $z$ axis:
% Source pp. 69 and 86 show the same two-rotation construction.
% Reused at both appearances of the source illustration.
\begin{center}
\begin{tikzpicture}[>=Latex,scale=1.08, every node/.style={font=\footnotesize}]
\coordinate (O) at (0,0);
\draw[->] (O)--(0,2.70) node[above] {$z$};
\draw[->] (O)--(2.55,-.48) node[right] {$y$};
\draw[->] (O)--(-1.80,-1.03) node[left] {$x$};
\draw[->,line width=0.9pt] (O)--(1.16,1.98)
  node[above right,xshift=1pt] {$\vect e_r$};
\draw[->,dashed,thick] (O)--(-.73,2.14)
  node[above left,xshift=-3pt] {$\vect e_r'$};
\draw[densely dashed] (0,1.72) ellipse [x radius=1.12,y radius=.30];
\draw[->] (-.12,1.36) arc[start angle=100,end angle=62,radius=.73];
\node at (-.70,1.07) {$\beta$};
\draw[->] (.55,1.98) arc[start angle=30,end angle=110,x radius=.50,y radius=.23];
\node at (.22,2.35) {$\alpha$};
\node[anchor=west] at (2.25,1.55) {second rotation about $z$};
\draw[->] (2.30,1.43) -- (1.20,1.55);
\node[anchor=west] at (2.25,.47) {first rotation about $y$};
\draw[->] (2.30,.38)--(1.16,.24);
\end{tikzpicture}
\end{center}
\[\ket{\vect e_r(\theta,\varphi)}
=\op D_z(\varphi)\op D_y(\theta)\ket{\vect e_z(0,\varphi)}.\]
Insert the complete basis:
\[\ket{\vect e_r(\theta,\varphi)}
=\sum_{m'}\op D_z(\varphi)\op D_y(\theta)\ket{l,m'}
\braket{l,m'}{\vect e_z(0,\varphi)}.\]
Hence
\begin{align*}
Y_{lm}^*(\theta,\varphi)
&=\braket{l,m}{\vect e_r(\theta,\varphi)}\\
&=\sum_{m'}\mel{l,m}{\op D_z(\varphi)\op D_y(\theta)}{l,m'}
\braket{l,m'}{\vect e_z(0,\varphi)}.
\end{align*}
Since
\[\braket{l,m'}{\vect e_z(0,\varphi)}=Y_{lm'}^*(0,\varphi)
=\sqrt{\frac{2l+1}{4\pi}}\delta_{m'0},\]
and
\[\bra{l,m}\op D_z(\varphi)
=\bra{l,m}e^{-\ii\op L_z\varphi/\hbar}
=e^{-\ii m\varphi}\bra{l,m},\]
we obtain
\[\boxed{Y_{lm}^*(\theta,\varphi)
=e^{-\ii m\varphi}\sqrt{\frac{2l+1}{4\pi}}\,d^l_{m0}(\theta),\qquad
 d^l_{mm'}(\theta)=\mel{l,m}{\op D_y(\theta)}{l,m'}.}\]
A spherical harmonic in any direction can thus be obtained by rotating a spherical harmonic along the $Z$ direction.